\section{The Product Tour}\label{sec:tour}

The three sections before this one built one machine: a single door, a handful of
checks, an append-only log, and views computed from it, fed by contracts that read
only declared data and stamped observations. The claim was that this one machine
carries every post-trade product, not a family of machines each shaped to its asset
class. A claim like that is not settled by argument. It is settled by running real
products through the machine and watching it hold.

That is what the next eight sections do. Each takes one product family, picks the
event on it that is hardest for the design, and works that event through in full.
The stops are ordered by machinery, not by asset class: the trade and the split
first, because they are the plainest triggers; then two lifecycle products, a listed
future and a structured note; then corporate actions in ascending order of
bespokeness, from a special-dividend override to a spin-off to an elective merger;
and securities lending last, because it is the one stop that changes what a position
\emph{is}, and so tests whether the architecture survives that change.

Every stop keeps the same five-part shape, so that what differs between products is
visible against what does not.

\begin{itemize}[leftmargin=1.5em]
\item \textbf{The trigger.} What fires --- an executed trade, a dated obligation
coming due, a confirmed corporate-action notice.
\item \textbf{The contract that catches it.} Which of the three trigger classes runs,
and what declared data and stamped observations it reads.
\item \textbf{The transactions it produces.} The moves and status writes it submits,
shown as numbered rows with their exact quantities.
\item \textbf{What the door checks.} Which admission check is load-bearing here, and
what it refuses.
\item \textbf{The design choice it substantiates.} One sentence, at the close, naming
the single decision from the mechanism sections this product would break without.
\end{itemize}

The stops share one setting --- a single firm's book over one year --- but each is a
self-contained worked example, and each opens with its own stated starting position
rather than inheriting a share count from the stop before. A recurring name, ACME,
carries the narrative: the same firm trades ACME, sees it split two-for-one, receives
a dividend forecast on it, holds an option over it through a September special
dividend, watches it spin off a subsidiary, and the following March sees it acquired
in an elective merger; the securities-lending stop works that same September through a
book that has lent part of its ACME. The tour makes no claim that one exact holding is
carried unchanged from stop to stop --- each stop states the position it works with, so
that what the machine does to that position is visible without the reader tracking a
running balance across a year. Two stops run their own lines, because the architecture
is not equity-shaped and the tour should not pretend it is: a dollar-denominated listed
future, and a structured note whose pricing the ledger never sees. Every number within
a stop is exact, and every one is a number the log would produce for the position that
stop states.
